\documentclass[11pt]{article}

\usepackage[margin=2.2cm]{geometry}
\usepackage{amsmath}
\usepackage{listings}
\usepackage{xcolor}
\usepackage{hyperref}
\usepackage{enumitem}

\hypersetup{colorlinks=true, linkcolor=blue!50!black, urlcolor=blue!50!black}

\lstdefinestyle{py}{
    backgroundcolor=\color{gray!5},
    basicstyle=\ttfamily\footnotesize,
    breaklines=true,
    frame=single,
    rulecolor=\color{gray!40},
    numbers=left,
    numberstyle=\tiny\color{gray!60},
    language=Python,
    tabsize=4,
    columns=fullflexible,
}
\lstset{style=py}

\setlength{\parskip}{0.4em}
\setlength{\parindent}{0pt}

\title{The Robot Adapter: Navigation Techniques in OVMM}
\author{}
\date{}

\begin{document}
\sloppy
\maketitle

\section*{Role of the Robot Adapter}

\texttt{core/missions/search\_and\_fetch.py} defines a \texttt{RobotAdapter} protocol
(\texttt{navigate\_to}, \texttt{navigate\_near}, \texttt{look\_at}, \texttt{get\_rgbd},
\texttt{grasp}) entirely in pure Python, with no ROS dependency. The mission logic --
``go to a candidate location, look, detect, grasp, repeat'' -- is written once against
that protocol. \texttt{ros2\_ws/src/fetcher/fetcher/robot\_adapter.py}
(\texttt{HsrRobotAdapter}) is the concrete implementation for the HSR: it is the only
place where Nav2, TF, camera topics, and joint trajectories actually appear. This seam
keeps the mission runner testable with a dry-run mock and reusable if the robot
platform ever changes.

\textbf{Frame handling.} The scene graph stores ground-truth poses in Gazebo's
\texttt{world} frame; Nav2 plans in the \texttt{map} frame, which is anchored at the
robot's spawn pose, not at \texttt{world} $(0,0)$. Every \texttt{navigate\_to}/
\texttt{navigate\_near}/\texttt{look\_at} call stamps its input as \texttt{world} and
lets a \texttt{tf2} buffer convert it, rather than assuming the two frames coincide.

\section*{Navigation: Reachability-Aware Approach Poses}

The first navigation approach picked a nav goal a fixed \texttt{standoff} (0.8\,m)
directly along $-X$ from the target's centroid, regardless of what was actually there.
This ignored both the environment (could land inside furniture, or facing a wall) and
the manipulator (a deep shelf could put the object well beyond the arm's working
range, producing visibly overstretched grasps). The fix replaces the fixed offset with
a \emph{reachability-aware} search against Nav2's live global costmap:

\begin{enumerate}[itemsep=2pt, topsep=2pt]
  \item Fetch the global costmap (\texttt{BasicNavigator.getGlobalCostmap()}) and
        convert it to a plain \texttt{CostmapGrid} (a small dataclass with no ROS
        message dependency, so the search logic itself stays pure Python and unit
        -testable: \texttt{core/perception/active\_perception.py}).
  \item Around the target $(c_x, c_y)$, scan an annulus of candidate standoff points
        $$ (v_x, v_y) = (c_x + r\cos\theta_i,\; c_y + r\sin\theta_i), \qquad
           r \in [\,r_{\min}, r_{\max}\,],\;\; \theta_i = \tfrac{2\pi i}{n} $$
        checking $r$ closest-first, so the tightest reachable spot wins over a
        farther one.
  \item Accept the first point whose costmap cost is below 253 (nav2\_costmap\_2d's
        inscribed-obstacle threshold) -- i.e.\ genuinely free, not merely
        ``discouraged.'' Its yaw faces back toward the target via $\mathrm{atan2}$,
        so the nav goal itself doesn't waste a step facing the wrong way.
\end{enumerate}

This says nothing about how precisely Nav2 actually achieves that goal, though --
real navigation has residual position/orientation error. That is handled separately,
described next.

\begin{lstlisting}[caption={Core search, \texttt{active\_perception.reachable\_approach\_pose}}]
for radius in np.linspace(reach_min, reach_max, n_radii):
    for i in range(n_angles):
        angle = 2 * math.pi * i / n_angles
        vx, vy = cx + radius * math.cos(angle), cy + radius * math.sin(angle)
        cost = costmap.cost_at(vx, vy)
        if cost is not None and cost < COSTMAP_BLOCKED_THRESHOLD:
            return _facing_pose(cx, cy, vx, vy)
\end{lstlisting}

$r_{\min}$/$r_{\max}$ are \emph{not} a fixed pair of constants. An early version tied
them to one global guess at the arm's comfortable reach, which doesn't generalize: a
band measured from a target's centroid is too close for furniture with real depth (an
HSR ended up base-to-edge against a table in practice) and needlessly far for a small
object. They are now computed per candidate in
\texttt{core/missions/search\_and\_fetch.py}'s \texttt{\_reach\_band()} as a fixed
clearance (\texttt{CLEARANCE\_MIN}/\texttt{CLEARANCE\_MAX}, roughly the robot's own
footprint radius) \emph{plus} that candidate's own \texttt{approach\_radius} --
half of the supporting furniture's \emph{shorter} horizontal dimension
(\texttt{core/reasoning/scene\_query.py}'s \texttt{\_approach\_radius()}), deliberately
the shorter one since an approach is normally roughly perpendicular to a piece's long
axis. A small side table and a long counter each get a standoff scaled to their own
size instead of one constant tuned for neither.

\textbf{Graceful degradation.} \texttt{navigate\_near} falls back to the original
fixed-standoff \texttt{approach\_pose()} if the world$\to$map TF lookup fails, the
costmap service is unavailable, or every cell in the reach band is blocked -- a
mission step never hard-fails just because the costmap-aware path couldn't run.

\section*{Active Perception, Re-centering, and the Grasp Handoff}

Once parked, \texttt{look\_at} points the head camera at the target using pure
trigonometry in \texttt{head\_geometry.py}: a pan delta from the bearing to the point
in the current pan-link frame, and a tilt angle from its elevation in the base frame,
both clamped to the joint limits read from the URDF. Both angles come from a live TF
lookup of the robot's actual current pose -- not the nav goal it was aiming for -- so
\texttt{look\_at} is already self-correcting for navigation error \emph{within} a
single call.

\textbf{Re-centering on the actual detection.} The first \texttt{look\_at} still aims
at the a-priori centroid (from the scene graph or the LLM), not at the object itself,
so if navigation stopped a little off from the planned approach pose, that first SAM3
detection can land off-center or near the edge of the frame.
\texttt{search\_and\_fetch.\_observe\_and\_detect} now treats that detection as
provisional: it lifts the mask to a point cloud, takes its centroid, converts that
camera-frame point back to world frame via a new adapter method
(\texttt{localize\_point}, a live \texttt{CAMERA\_FRAME}$\to$\texttt{WORLD\_FRAME} TF
lookup -- the mirror image of \texttt{look\_at}'s own conversion), and calls
\texttt{look\_at} a second time on \emph{that} -- the object's real position, not the
estimate. SAM3 runs once more on the re-centered view, and that second detection is
what gets handed to \texttt{grasp()}. Either TF lookup or re-detection failing just
falls back to the original detection rather than losing the object:
\begin{lstlisting}[caption={search\_and\_fetch.py -- re-centering before grasping}]
points = np.asarray(mask_to_pointcloud(mask, depth, K).points)
center_world = robot.localize_point(points.mean(axis=0).tolist())
if center_world is None:
    return None                  # caller falls back to the original detection
robot.look_at(center_world)      # aim at the object itself, not the a-priori estimate
rgb2, depth2, K2 = robot.get_rgbd()
det2 = sam3.detect(rgb2, target)
\end{lstlisting}

After that, \texttt{grasp()} hands the (now centered) point cloud to AnyGrasp,
transforms the chosen grasp pose from camera frame to \texttt{base\_footprint} via a
TF lookup taken at that exact moment, and converts it to arm joints with the same kind
of analytic formula (no IK solver) -- AnyGrasp's own candidate ranking already
accounts for collisions, so the adapter only needs to reach the pose it is given.

\section*{Takeaway}

The adapter's real job is reconciling three coordinate conventions (world ground
truth, Nav2 map, robot base/camera) and never fully trusting where it \emph{meant} to
end up: the costmap search turns an approximate target location into a reachable,
collision-checked one instead of a hardcoded offset, and re-centering on the actual
detection corrects for whatever residual error navigation still left, by looking
again rather than assuming the first view was good enough. Both are the same
principle -- consult live data instead of trusting the plan -- and it generalizes to
any future grasp-region or approach-direction refinement.

\end{document}
